
\documentclass{article}
\usepackage{colortbl}
\usepackage{makecell}
\usepackage{multirow}
\usepackage{supertabular}

\begin{document}

\newcounter{utterance}

\centering \large Interaction Transcript for game `cladder', experiment `full\_v1.5\_default', episode 6126 with selfrefine{-}llama{-}3.1{-}8b{-}instruct.
\vspace{24pt}

{ \footnotesize  \setcounter{utterance}{1}
\setlength{\tabcolsep}{0pt}
\begin{supertabular}{c@{$\;$}|p{.15\linewidth}@{}p{.15\linewidth}p{.15\linewidth}p{.15\linewidth}p{.15\linewidth}p{.15\linewidth}}
    \# & \multicolumn{2}{c}{Player} && \multicolumn{2}{c}{Game Master} \\
    \hline

    \theutterance \stepcounter{utterance}  
    & & & \multicolumn{4}{p{0.6\linewidth}}{
        \cellcolor[rgb]{0.9,0.9,0.9}{
            \makecell[{{p{\linewidth}}}]{
                \texttt{\tiny{[P1$\langle$GM]}}
                \texttt{You are an expert in causal inference. The following question is not a typical commonsense query, but rather a meticulously designed question created by a professor specializing in causal inference, intended to assess the students' mastery of the course content.} \\
\texttt{Imagine a self{-}contained, hypothetical world with only the following conditions, and without any unmentioned factors or causal relationships: Smoking has a direct effect on effort and lactose intolerance. Effort has a direct effect on lactose intolerance. We know that smoking causes being lazy. smoking and being hard{-}working causes being lactose intolerant. Would the student is lactose intolerant if nonsmoking instead of smoking?} \\
\texttt{Start your answer with 'Yes' or 'No', followed by additional reasoning or evidence to support your explanation.} \\
            }
        }
    }
     \\ \\

    \theutterance \stepcounter{utterance}  
    & \multicolumn{4}{p{0.6\linewidth}}{
        \cellcolor[rgb]{0.561,0.737,0.561}{
            \makecell[{{p{\linewidth}}}]{
                \texttt{\tiny{[P1$\rangle$GM]}}
                \texttt{No.} \\
\\ 
\texttt{To address this question, we need to apply the principles of causal inference, specifically the concept of counterfactuals and the backdoor criterion. } \\
\\ 
\texttt{First, let's establish the causal relationships:} \\
\\ 
\texttt{{-} Smoking causes being lazy (S → L)} \\
\texttt{{-} Smoking causes being lactose intolerant (S → LI)} \\
\texttt{{-} Being hard{-}working causes being lactose intolerant (HW → LI)} \\
\texttt{{-} Being lazy causes being hard{-}working (L → HW)} \\
\\ 
\texttt{We are asked to consider the counterfactual scenario where the student is a nonsmoker instead of a smoker (S' = 0). We want to know if the student would be lactose intolerant in this scenario (LI').} \\
\\ 
\texttt{To answer this question, we need to consider the backdoor paths from smoking to lactose intolerance. There are two backdoor paths:} \\
\\ 
\texttt{1. S → L → HW → LI} \\
\texttt{2. S → LI} \\
\\ 
\texttt{To block these backdoor paths, we need to condition on the variables that are on the backdoor paths. In this case, we need to condition on being lazy (L) and being hard{-}working (HW).} \\
\\ 
\texttt{However, conditioning on being lazy (L) would not block the backdoor path S → L → HW → LI, because being lazy is a direct effect of smoking. Therefore, we cannot condition on being lazy (L).} \\
\\ 
\texttt{On the other hand, conditioning on being hard{-}working (HW) would block the backdoor path S → L → HW → LI, because being hard{-}working is a direct effect of being lazy, which is a direct effect of smoking. However, conditioning on being hard{-}working (HW) would also block the direct effect of smoking on lactose intolerance (S → LI), because being hard{-}working is a direct cause of lactose intolerance.} \\
\\ 
\texttt{Therefore, conditioning on being hard{-}working (HW) would not allow us to identify the effect of smoking on lactose intolerance. In other words, we cannot use the backdoor criterion to identify the counterfactual effect of smoking on lactose intolerance.} \\
\\ 
\texttt{However, we can use the frontdoor criterion to identify the effect of smoking on lactose intolerance. The frontdoor criterion states that if we have a direct effect of a variable on the outcome, and we can condition on the variables that are on the causal path, then we can identify the effect of the variable on the outcome.} \\
\\ 
\texttt{In this case, we have a direct effect of smoking on lactose intolerance (S → LI), and we can condition on being hard{-}working (HW). Therefore, we can use the frontdoor criterion to identify the effect of smoking on lactose intolerance.} \\
\\ 
\texttt{However, the question asks us to consider the counterfactual scenario where the student is a nonsmoker instead of a smoker. In this scenario, the direct effect of smoking on lactose intolerance (S → LI) is absent. Therefore, we cannot use the frontdoor criterion to identify the effect of smoking on lactose intolerance in the counterfactual scenario.} \\
\\ 
\texttt{In summary} \\
            }
        }
    }
    & & \\ \\

    \theutterance \stepcounter{utterance}  
    & & & \multicolumn{2}{p{0.3\linewidth}}{
        \cellcolor[rgb]{0.9,0.9,0.9}{
            \makecell[{{p{\linewidth}}}]{
                \texttt{\tiny{[GM$|$GM]}}
                \texttt{No} \\
            }
        }
    }
    & & \\ \\

    \theutterance \stepcounter{utterance}  
    & & & \multicolumn{2}{p{0.3\linewidth}}{
        \cellcolor[rgb]{0.9,0.9,0.9}{
            \makecell[{{p{\linewidth}}}]{
                \texttt{\tiny{[GM$|$GM]}}
                \texttt{no} \\
            }
        }
    }
    & & \\ \\

    \theutterance \stepcounter{utterance}  
    & & & \multicolumn{2}{p{0.3\linewidth}}{
        \cellcolor[rgb]{0.9,0.9,0.9}{
            \makecell[{{p{\linewidth}}}]{
                \texttt{\tiny{[GM$|$GM]}}
                \texttt{game\_result = WIN} \\
            }
        }
    }
    & & \\ \\

\end{supertabular}
}

\end{document}
